%lezione 21 dell'11 maggio 2026 pag 47 appunti
%-----------------------------------------------------------------------------------------------------

L'insieme $\left\{\varphi^0_{klm}(r,\vartheta,\varphi)\right\}$ è un sistema completo di
$\mathcal{H}=L^2(\mathbb{R}^3)$: per proseguire con la soluzione del problema di scattering è
necessario decomporre la soluzione omogenea $\psi_\text{hom}$ in questa base
$$\psi_\text{hom}(r,\vartheta,\varphi) = \frac{1}{(2\pi)^{\frac32}} e^{ikr\cos\vartheta} = 
\bra{\mathbf{x}}\ket{0,0,k}$$
dove $\ket{0,0,k} = \ket{l=0,m=0,k=k}$ è un autostato simultaneo di $L^2$ ($l=0$), $L_z$ ($m=0$) e
$H_0$ ($k=k$): è banale mostrare che $\psi_\text{hom}$ è la rappresentazione nello spazio delle
coordinate di tale autostato. 

In questa notazione, le $\varphi^0_{klm}$ rappresentano lo stato $\ket{l,m,k}$: utilizzando la
completezza di questi autostati, si trova che
$$\ket{0,0,k}=\int_{}^{}dk'\sum_{l'=0}^{\infty}\sum_{m'=-l'}^{l'}\ket{l',m',k'}\bra{l',m',k'}
\ket{0,0,k} =\sum_{l'=0}^{\infty}\bra{l',0,k}\ket{0,0,k}\ket{l',0,k}
:=\sum_{l'=0}^{\infty}c_{kl'}\ket{l',0,k}$$
proiettando tale espressione nello spazio delle coordinate si ha
$$\bra{\mathbf{x}}\ket{0,0,k} = \sum_{l'=0}^{\infty}c_{kl'}\bra{\mathbf{x}}\ket{l',0,k} \rightarrow
\frac{e^{ikz}}{(2\pi)^\frac32}= \sum_{l'=0}^{\infty}c_{kl'}\sqrt{\frac{2k^2}{\pi}}j_{l'}(kr)
Y^{l'0}(\vartheta,\varphi)
$$
Le armoniche sferiche con $m=0$ sono proporzionali ai polinomi di Legendre di argomento
$\cos\vartheta$: assorbendo tutte le costanti in $d_{kl'}$ si trova
$$e^{ikr\cos\vartheta}=\sum_{l'=0}^{\infty}d_{kl'}j_{l'}(kr)P_{l'}(\cos\vartheta)$$
Si vogliono ora determinare i coefficienti $d_{kl'}$: si ridefinisce dunque $x:=\cos\vartheta$,
si moltiplicano entrambi i membri dell'espressione per $P_l(x)$ e si integra sul dominio
$x\in[-1,1]$
$$\int_{-1}^{1}dxe^{ikrx}P_l(x) = \sum_{l'=0}^{\infty}d_{kl'}j_{l'}(kr)\int_{-1}^{1}dxP_l(x)P_{l'}(x)
=\frac{2}{2l+1}d_{kl}j_l(kr)
$$
dove si è utilizzata l'ortogonalità dei $P_l$
$$\int_{-1}^{1}dxP_l(x)P_{l'}(x)=\frac{2}{2l+1}\delta_{ll'}$$
Si è giunti all'espressione
\begin{equation}
    I:=\int_{-1}^{1}dxe^{ikrx}P_l(x) = \frac{2d_{kl}}{2l+1}j_l(kr)
    \label{eq:dkl}
\end{equation}
Ricordando la formula di Rodriguez per i polinomi di Legendre
$$P_l(x) = \frac{(-1)^l}{2^ll!} \dv[l]{}{x}\left(1-x^2\right)^l$$
si calcola ora il lato sinistro della (\ref{eq:dkl})
$$I = \int_{-1}^{1}e^{ikrx}P_l(x) = 
\frac{(-1)^l}{2^ll!}\int_{-1}^{1}dx \dv[l]{}{x} \left(1-x^2\right)^le^{i\rho x}
$$
integrando per parti $l$ volte, i termini di bordo sono $\eval{\propto(1-x^2)}_{-1}^1=0$ e
si ha
$$I = \frac{1}{2^ll!}\int_{-1}^{1}dx (1-x^2)^l(i\rho)^le^{i\rho x}$$
Poiché i coefficienti $d_{kl}$ sono indipendenti da $\rho$, si considera il limite
$\rho\to0$ e si trova
$$I\underset{\rho\to0}\sim\frac{1}{2^ll!}\int_{-1}^{1}dx (1-x^2)^l(i\rho)^l(1+O(\rho)) = 
\frac{i^l\rho^l}{2^ll!}\int_{-1}^{1}dx (1-x^2)^l = \frac{i^l\rho^l}{2^ll!}B(l+1,l+1)2^{2l+1}
$$
dove $B(l+1,l+1)$ è la Beta di Eulero: una delle proprietà di questa funzione è
$$B(u,v) = \frac{\Gamma(u)\Gamma(v)}{\Gamma(u+v)}$$
dunque si ha
$$I\underset{\rho\to0}\sim\frac{i^l\rho^l}{2^ll!}\frac{l!l!}{(2l+1)!}2^{2l+1} = 2i^l\rho^l
\frac{2^ll!}{(2l+1)!} = \frac{2i^l\rho^l}{(2l+1)!!}$$
dove si è utilizzato il fatto che $2^ll!=(2l)!!$ e $\frac{(2l)!!}{(2l+1)!} = \frac{1}{(2l+1)!!}$.

Il lato destro della (\ref{eq:dkl}) per $\rho\to0$ vale
$$\frac{2d_{lk}}{2l+1}j_l(\rho)\underset{\rho\to0}\sim\frac{2}{2l+1}d_{kl}\frac{\rho^l}{(2l+1)!!}$$
eguagliando le due espressioni si trova che
$$\rho\to0 \longrightarrow \frac{2i^l\rho^l}{(2l+1)!!} = \frac{2}{2l+1}d_{kl}\frac{\rho^l}{(2l+1)!!}
\leftrightarrow d_{kl} = (2l+1)i^l
$$
Si è dunque trovato che
$$\boxed{\psi_\text{hom}(\mathbf{x})=e^{ikz}=\sum_{l=0}^{\infty}i^l(2l+1)j_l(kr)P_l(\cos\vartheta)}$$
\subsection{Soluzione con il potenziale}
Conclusa la scomposizione in onde parziali della particella libera, si considera ora la funzione
d'onda soggetta al potenziale $V(r)$: per la simmetria del problema rispetto a SO(3) (potenziale
centrale), ci si aspetta una funzione d'onda e un'ampiezza differenziale indipendente da $\varphi$
\begin{equation}
\psi_k(\mathbf{x}) = \psi_k(r,\vartheta) = e^{ikr\cos\vartheta} + f_k(\vartheta)\frac{e^{ikr}}{r}
\label{eq:sol_pot}
\end{equation}

Come prima, si scompone la funzione d'onda in parte radiale e angolare e, per questa simmetria, la
parte angolare sarà indipendente da $\varphi$:
\begin{equation*}
\psi_k(r,\vartheta) = \sum_{l=0}^{\infty}i^l(2l+1)\varphi_l(kr)P_l(\cos\vartheta)
\end{equation*}

\textit{Osservazione} La forma della soluzione è la medesima di quella della particella libera, con la
differenza che $\varphi_l(kr)$ non è semplicemente la $j_l(kr)$, ma è la soluzione dell'equazione
radiale con il potenziale.

L'equazione radiale è
$$\varphi''_l(\rho) + \frac{2}{\rho}\varphi'_l(\rho) + \left(1-\frac{l(l+1)}{\rho^2}+U(\rho)\right)
\varphi_l(\rho)=0$$
Per ipotesi $U(\rho) = 0\phantom{x}\forall \rho>kA$, dunque per $r>A$
l'equazione si riduce alla Bessel sferica, con soluzione
$$\varphi_l(\rho)\underset{r>A}=\alpha_lj_l(\rho)+\beta_l\eta_l(\rho)$$
Poiché $U(\rho)\in\mathbb{R}$, se $\varphi_l$ è soluzione, lo è anche $\varphi_l^*$, ma le
$j_l,\eta_l$ sono funzioni reali, dunque $\varphi^*_l$ non è una soluzione linearmente indipendente:
$$\varphi_l=A\varphi^*_l \leftrightarrow
\begin{dcases}
\alpha^*_l = A\alpha_l \\
\beta^*_l= A\beta_l
\end{dcases}\leftrightarrow
\left(\frac{\alpha_l}{\beta_l}\right)^* = \frac{\alpha_l}{\beta_l}\in\mathbb{R}\rightarrow
\frac{\beta_l}{\alpha_l}:=\tan(\delta_l)
$$
dunque
$$\varphi_l(kr) \underset{r>A}= \alpha_l\left(j_l(kr)+\tan(\delta_l)\eta_l(kr)\right) = 
A_l(\cos(\delta_l)j_l(kr)+\sin(\delta_l)\eta_l(kr))
$$
con $A_l:=\flatfrac{\alpha_l}{\cos\delta_l}$

Si consideri ora l'andamento di quest'ultima espressione per $r\to\infty$:
\begin{gather*}
\varphi_l(kr)\underset{r\to\infty}\sim\frac{A_l}{kr}\left(\cos(\delta_l)\sin\left(kr-\frac{l\pi}{2}
\right)+\sin(\delta_l)\cos\left(kr-\frac{l\pi}{2}\right)\right)=\frac{A_l}{kr}
\sin\left(kr-\frac{l\pi}{2}+\delta_l\right)=\\
=\frac{A_l}{2ikr}\left[e^{ikr}e^{-i\frac{l\pi}{2}}e^{i\delta_l}-
e^{-ikr}e^{i\frac{l\pi}{2}}e^{-i\delta_l}\right]
\end{gather*}
\textit{Osservazione} Come per la particella libera, anche questa soluzione è la sovrapposizione di
un'onda entrante con un'onda uscente, che viene riflessa tanto più lontano dal centro diffusore quanto
$l$ è grande: con un ragionamento analogo al precedente, per basse energie i valori di $l$ che
contribuiscono sono pochi ed è proprio questo il regime in cui si usa questa approssimazione;
diversamente da prima, lo sfasamento qui vale $l\pi+2\delta_l$: da questo sfasamento, che, a monte,
dipende dal fatto che in questo caso non si è scartata la soluzione $\eta_l$ della Bessel sferica,
si troveranno le informazioni sull'interazione avvenuta.

L'andamento della soluzione sarà dunque
$$\psi_k(r,\vartheta) = \sum_{l=0}^{\infty}i^l(2l+1)\varphi_l(kr)P_l(\cos\vartheta)
\underset{r\to\infty}\sim
\sum_{l=0}^{\infty}A_li^l(2l+1)\left[e^{ikr}e^{-i\frac{l\pi}{2}}e^{i\delta_l}-
e^{-ikr}e^{i\frac{l\pi}{2}}e^{-i\delta_l}\right]P_l(\cos\vartheta)
$$
Questo andamento va confrontato con quello della (\ref{eq:sol_pot}), che è
$$\psi_k(r,\vartheta) = e^{ik\cos\vartheta}+f(\vartheta)\frac{e^{ikr}}{r}\underset{r\to\infty}\sim
\sum_{l=0}^{\infty}i^l(2l+1)\frac{1}{2ikr}\left[e^{ikr}e^{-i\frac{l\pi}{2}}-
e^{-ikr}e^{i\frac{l\pi}{2}}\right]P_l(\cos\vartheta)+f(\vartheta)\frac{e^{ikr}}{r}
$$
Si confrontano i coefficienti di $\flatfrac{e^{-ikr}}{r}$:
$$\frac{e^{-ikr}}{2ikr}\sum_{l=0}^{\infty}i^l(2l+1)
e^{i\frac{l\pi}{2}} = \frac{e^{-ikr}}{2ikr}\sum_{l=0}^{\infty}i^l(2l+1)A_le^{l\frac{\pi}{2}}
e^{-i\delta_l}\leftrightarrow A_l=e^{i\delta_l}$$
e quelli di $\flatfrac{e^{ikr}}{r}$

\begin{gather*}
\frac{e^{ikr}}{r}\left[f_k(\vartheta) +
\frac{1}{k}\sum_{l=0}^{\infty}i^l(2l+1)\frac{1}{2i}e^{-i\frac{l\pi}{2}}P_l\right] = 
\frac{e^{ikr}}{r}\left[\frac{1}{2ik}\sum_{l=0}^{\infty}i^l(2l+1)e^{-i\frac{l\pi}{2}}
e^{i\delta_l}P_l\right] \leftrightarrow \\
\boxed{
f(\vartheta)=\frac{1}{2ik}\sum_{l=0}^{\infty}(2l+1)e^{i\delta_l}\sin(\delta_l)P_l(\cos\vartheta)
}\end{gather*}
Da questa espressione si calcolano infine la sezione d'urto differenziale
$$\dv[]{\sigma}{\Omega}(\vartheta) = |f_k(\vartheta)|^2 = \frac{1}{k^2}
\sum_{l=0}^{\infty}\sum_{l'=0}^{\infty}(2l+1)(2l'+1)e^{i(\delta_l-\delta_{l'}}\sin(\delta_l)
\sin(\delta_{l'})P_l(\cos\vartheta)P_{l'}(\cos\vartheta)
$$
e quella totale
$$\sigma_\text{tot} = \int_{}^{}d\Omega\dv[]{\sigma}{\Omega} = 2\pi\int_{-1}^{1}d\cos\vartheta \frac{1}
{k^2} \sum_{l,l'=0}^{\infty}[...]P_l(\cos\vartheta)P_{l'}(\cos\vartheta) = \frac{4\pi}{k^2}(2l+1)\sin^2
(\delta_l)$$
\textit{Osservazione 1} La sezione d'urto differenziale presenta dei termini di "interferenza" tra le
onde parziali in cui si è scomposto il problema, mentre quella totale no: questo è dovuto al fatto che
le $\varphi_l$ hanno una dipendenza da $\vartheta$ e dunque interferiscono a $\vartheta$ fissato,
mentre se si considera tutto l'angolo solido questa dipendenza sparisce.

\textit{Osservazione 2} Il teorema ottico è soddisfatto, in quanto
$$f_k(0) = \frac{1}{k}\sum_{l=0}^{\infty} (2l+1)e^{i\delta_l}\sin(\delta_l)\rightarrow
\frac{4\pi}{k}\Im(f_k(0)) = \frac{4\pi}{k^2}(2l+1)\sin^2(\delta_l)=\sigma_\text{tot}
$$
\subsection{Esempio: la sfera impenetrabile}
La sfera impenetrabile è un potenziale centrale definito come
$$V(r)=\begin{dcases}
    +\infty &\text{se }x<A\\
    0 &\text{se }x>A
\end{dcases}
$$
si determina prima la sezione d'urto geometrica: una particella che urta una sfera rigida con
parametro d'impatto $b = A\cos\left(\frac{\vartheta}{2}\right)$, dove $\vartheta$ è l'angolo di
deflessione della particella rispetto all'asse $z$; si considera ora la corona circolare definita da
$b+db$: la sua area vale
$$d\sigma = 2\pi b|db| = 2\pi A\cos\left(\frac{\vartheta}{2}\right)A\left|\left(
-\frac{1}{2}\sin\left(\frac{\vartheta}{2}\right)d\vartheta\right)\right|$$
L'angolo solido che essa sottende è invece
$$d\Omega = 2\pi\sin\vartheta d\vartheta = 4\pi\sin\left(\frac{\vartheta}{2}\right)
\cos\left(\frac{\vartheta}{2}\right)d\vartheta$$
La sezione d'urto differenziale geometrica vale dunque
$$\dv[]{\sigma}{\Omega} = \frac{A^2\pi\cos\left(\frac{\vartheta}{2}\right)
\sin\left(\frac{\vartheta}{2}\right)d\vartheta}{4\pi\cos\left(\frac{\vartheta}{2}\right)
\sin\left(\frac{\vartheta}{2}\right)d\vartheta} = \frac{A^2}{4}$$
Mentre quella totale
$$\sigma_\text{tot} = \int_{}^{}d\Omega\dv[]{\sigma}{\Omega} = \pi A^2$$
che è, come ci si aspetta, l'area del cerchio massimo della sfera.

Utilizzando invece la decomposizione in onde parziali, si trova che le $\varphi_l$ valgono
$$\varphi_l(kr) = \begin{dcases}
    0 &\text{se } r<A\\
    A_l(\cos(\delta_l)j_l(kr) + \sin(\delta_l)\eta_l(kr)) &\text{se }r>A
\end{dcases}
$$
Si impone la continuità in $r=A$:
$$\tan(\delta_l) = \frac{j_l(kA)}{\eta_l(kA)}$$
Per i risultati precedenti, la sezione d'urto totale vale
$$\sigma_\text{tot} = \frac{4\pi}{k^2}\sum_{l=0}^{\infty}(2l+1)\frac{\tan^2(\delta_l)}
{1-\tan^2(\delta_l)}=\frac{4\pi}{k^2}\sum_{l=0}^{\infty}(2l+1)\frac{j^2_l(kA)}{j_l^2(kA)+\eta_l^2(kA)}
$$
Nel regime di basse energie, $kA<<1\rightarrow\delta_l<<1$ e, utilizzando gli andamenti asintotici
della $j_l$ e della $\eta_l$, si trova
$$\tan(\delta_l)\sim\delta_l \sim -\frac{(kA)^l}{(2l+1)!!}\frac{(kA)^{l+1}(2l+1)}{(2l+1)!!} = 
\frac{-(2l+1)}{((2l+1)!!)^2}(kA)^{2l+1}
$$
È dunque evidente che in questo regime il valore di $\delta_l$ decresce molto velocemente per $l$
crescenti: per calcolare la sezione d'urto si considera dunque solo il valore $l=0$
$$\delta_0 = \frac{j_0(kA)}{\eta_0(kA)} = -\tan(kA)\sim-kA$$
dunque
$$f_k(\vartheta) = \frac{1}{k}(2\cdot0+1)e^{i\delta_0}\sin(\delta_0) = -\frac{1}{k}e^{-ikA}\sin(kA)$$
e infine
$$\sigma_\text{tot} = \frac{4\pi}{k^2}\sin^2(kA)\underset{kA\to0}\sim 4\pi A^2$$
\textit{Osservazione 1} La sezione d'urto differenziale è isotropa perché, nell'approssimazione di
$kA$ piccolo, la particella incidente ha $\lambda$ Compton grande e non "vede" la sfera.

\textit{Osservazione 2} La sezione d'urto totale è 4 volte quella geometrica per la natura ondulatoria
della particella incidente.
\newpage
